% ============================================================
%  Krimelack: Content-Addressable Structural Memory
%  with Clockless Habituation for Ternary Perception Systems
%  Version 1.0 — May 2026
%  Confidential — Internal Use Only
% ============================================================
\documentclass[11pt, a4paper]{article}

\usepackage[margin=2.2cm]{geometry}
\usepackage{amsmath, amssymb, amsthm, mathtools}
\usepackage{booktabs, array, longtable}
\usepackage{xcolor}
\usepackage{hyperref}
\usepackage{titlesec}
\usepackage{enumitem}
\usepackage{fancyhdr}
\usepackage{setspace}
% algorithm packages removed (not available in this TeX distribution)

\hypersetup{
  colorlinks=true,
  linkcolor=blue!50!black,
  urlcolor=blue!60!black,
}

\definecolor{nm-dark}{RGB}{30, 30, 30}
\definecolor{nm-gray}{RGB}{100, 100, 100}
\definecolor{nm-green}{RGB}{20, 100, 40}

\pagestyle{fancy}
\fancyhf{}
\rhead{\textcolor{nm-gray}{\small Krimelack Paper v1.0}}
\lhead{\textcolor{nm-gray}{\small Confidential --- Internal Use Only}}
\cfoot{\thepage}

\titleformat{\section}{\large\bfseries\color{nm-dark}}{\thesection}{0.6em}{}[\titlerule]
\titleformat{\subsection}{\normalsize\bfseries\color{nm-dark}}{\thesubsection}{0.6em}{}
\titleformat{\subsubsection}{\normalsize\itshape\color{nm-dark}}{\thesubsubsection}{0.6em}{}

\setlength{\parindent}{0pt}
\setlength{\parskip}{6pt}
\onehalfspacing

\newtheorem{definition}{Definition}[section]
\newtheorem{theorem}{Theorem}[section]
\newtheorem{proposition}{Proposition}[section]
\newtheorem{remark}{Remark}[section]

\begin{document}

\begin{center}
  {\LARGE \textbf{Krimelack}}\\[6pt]
  {\large Content-Addressable Structural Memory with}\\[2pt]
  {\large Clockless Habituation for Ternary Perception Systems}\\[10pt]
  {\color{nm-gray}\small Version 1.0 \quad|\quad May 2026 \quad|\quad Confidential}\\[4pt]
  {\color{nm-gray}\small Joseph Forrester \quad|\quad DSF-AI Architecture}
\end{center}

\vspace{1em}
\hrule
\vspace{1em}

% ============================================================
\section*{Abstract}

We present Krimelack, a content-addressable structural memory
architecture for balanced ternary perception systems. Krimelack
stores \textit{structural motifs} --- settled states of a
combinational coupling fabric --- and provides cue-based recall
with clockless habituation feedback. Unlike conventional
content-addressable memories (CAMs) that match bit patterns,
Krimelack stores coupled multi-field patterns in balanced ternary
($-1, 0, +1$) and scores matches by counting agreeing non-null
trits, giving the null state ($0$) a computational role: it
represents structural uncertainty, not a value to be matched.

The architecture combines four capabilities that do not exist
together in any prior system: (i)~structural motif storage with
quality-gated commits, (ii)~parallel cue-based recall from partial
signatures, (iii)~pattern-level habituation via combinational
dead-zone feedback, and (iv)~autonomous commit gating based on
dimensional exhaustion detection. Krimelack has been implemented
in Verilog, validated on Xilinx Zynq-7020 FPGA silicon, and
demonstrated in autonomous robot navigation with camera and IR
sensor fusion (351 samples, 176 seconds continuous operation).

\textbf{Gate count:} $\sim$4,600 equivalent gates per block
(64-entry, 48-bit). \textbf{Recall latency:} combinational
(zero clock cycles). \textbf{Prior art:} none for the combined
architecture.

\tableofcontents

% ============================================================
\section{Introduction}

\subsection{The Problem}

Perception systems that operate without training --- where a
single exposure must suffice for future recognition --- require
a memory architecture fundamentally different from neural network
weight matrices or conventional lookup tables.

The requirements are:
\begin{enumerate}[nosep]
  \item \textbf{Store structural patterns, not raw data.}
        The memory must capture the settled state of a
        perception system, not the sensor input that caused it.
  \item \textbf{Recall from partial cues.}
        A noisy or incomplete sensor input must retrieve the
        closest stored pattern, not fail on exact mismatch.
  \item \textbf{Reject noise.}
        Patterns captured during ambiguous or uncertain
        perception must be blocked from storage.
  \item \textbf{Habituate to familiarity.}
        Repeated exposure to the same environment must not
        produce the same response indefinitely. The system
        must seek novelty.
  \item \textbf{Operate without a clock in the recall path.}
        For integration with combinational perception
        hardware (nanosecond settling), recall must be
        combinational.
\end{enumerate}

No existing memory architecture satisfies all five requirements.

\subsection{Context: The SPPU Perception Fabric}

Krimelack is designed for the ArcLoom Structural Perception
Processing Unit (SPPU), a combinational balanced ternary
coupling fabric that settles to perception decisions through
wire propagation in $\sim$5~ns. The SPPU produces a
\textit{loom state}: a vector of balanced ternary trits
($-1, 0, +1$) encoded as 2-bit pairs across multiple
coupled strands.

The loom state captures what the perception system
\textit{committed to} given a set of sensor inputs. Trits
at $\pm 1$ represent structural commitments. Trits at $0$
represent dimensions where coupling pressure was insufficient
to force a decision --- structural uncertainty, not missing
data.

Krimelack stores these loom states as \textit{motifs} and
provides cue-based recall to enable temporal pattern recognition
across perception events.

\subsection{Naming}

``Krimelack'' derives from the concept of a ``crime of the
lacquer'' --- a structural record that cannot be erased once
committed, like a flaw sealed under a coat of lacquer. A
committed motif is a permanent structural fact, not a mutable
variable.

% ============================================================
\section{Architecture}

\subsection{Memory Structure}

\begin{definition}[Motif]
A \textbf{motif} is a stored loom state: a vector of $N$
balanced ternary trits encoded as $2N$ bits.
\begin{itemize}[nosep]
  \item \texttt{00} = null ($0$): structural uncertainty
  \item \texttt{01} = positive ($+1$): excitatory commitment
  \item \texttt{10} = negative ($-1$): inhibitory commitment
  \item \texttt{11} = invalid (never stored)
\end{itemize}
\end{definition}

The current FPGA implementation uses:
\begin{center}
\begin{tabular}{@{}lr@{}}
\toprule
\textbf{Parameter} & \textbf{Value} \\
\midrule
Trits per motif & 24 (8 strands $\times$ 3 settling/decision trits) \\
Bits per motif & 48 \\
Motif depth & 32 entries (circular buffer) \\
Hash width & 8 bits (XOR-fold duplicate detection) \\
\bottomrule
\end{tabular}
\end{center}

\subsection{Functional Blocks}

Krimelack consists of four functional blocks:

\begin{enumerate}
  \item \textbf{Motif Store:} A register array holding $D$ motifs
        of $W$ bits each, with a circular write pointer and
        occupancy counter.

  \item \textbf{Commit Gate:} Quality filter that determines
        whether a candidate motif is eligible for storage.
        Prevents committing noise, duplicates, or ambiguous
        states.

  \item \textbf{Recall Engine:} Parallel comparator that scores
        every stored motif against a query pattern simultaneously.
        Returns the best match and its score. Purely combinational
        --- zero clock cycles for recall.

  \item \textbf{Habituation Feedback:} The match score feeds back
        into the SPPU's dead-zone thresholds, raising the
        commitment barrier for familiar patterns. This is the
        mechanism by which Krimelack produces novelty-seeking
        behavior without any clock or timer.
\end{enumerate}

% ============================================================
\section{Commit Gating}

\subsection{The Problem of Noisy Commits}

A perception system that stores every loom state would fill its
memory with noise --- transient states, sensor glitches, and
ambiguous readings. Krimelack uses \textit{commit gating} to
ensure that only structurally significant states are stored.

\subsection{Commit Eligibility Criteria}

A motif is committed if and only if ALL of the following hold:

\begin{enumerate}[nosep]
  \item \textbf{Structural lock (L6):} The dimensional exhaustion
        detector has fired, indicating that the loom state has
        collapsed to a geometrically inevitable configuration.
        Formally: the number of non-null trits exceeds the
        threshold $n_{start}/e$, where $n_{start}$ is the total
        number of trits and $e = 2.718\ldots$ is Euler's number.

  \item \textbf{Low uncertainty:} The adjusted uncertainty
        $U^* \leq U_{settle}$ (default: $0.35$ in Q16.16
        fixed-point). This ensures the perception system is
        confident about the current state.

  \item \textbf{Sufficient resonance:} The resonance score
        $R \geq R_{commit}$ (default: $0.72$ in Q16.16).
        This ensures the state has structural coherence, not
        just low noise.

  \item \textbf{SafeMode inactive:} The system is not in a
        protective mode that suppresses storage.

  \item \textbf{No duplicate:} The candidate motif does not
        already exist in memory. Checked via hash-based
        comparison (XOR-fold to 8 bits, followed by exact
        bit comparison on hash collision).
\end{enumerate}

\subsection{Forced Commit (Software Capture)}

For calibration and training scenarios (e.g., turntable object
capture), software can trigger a commit via AXI register write.
The forced commit bypasses criteria (1)--(4) but still enforces
criterion (5) --- duplicate detection is never disabled. This
ensures idempotency: presenting the same stimulus twice never
produces two motifs.

\subsection{Formal Specification}

\begin{equation}
  \text{commit} =
  \begin{cases}
    1 & \text{if } (\text{L6\_fire} \wedge U^* \leq U_s
        \wedge R \geq R_c \wedge \neg\text{safe}
        \wedge \neg\text{dup}) \\
      & \quad \vee\; (\text{force} \wedge \neg\text{dup}) \\
    0 & \text{otherwise}
  \end{cases}
\end{equation}

% ============================================================
\section{Recall: Parallel Pattern Matching}

\subsection{Scoring Function}

The recall engine compares a query pattern $Q$ against every
stored motif $M_k$ ($k = 0, \ldots, D-1$) simultaneously.
The score counts the number of trits where both the query
and the motif are non-null and agree:

\begin{equation}
  \text{score}(Q, M_k) = \sum_{i=0}^{N-1}
  \begin{cases}
    1 & \text{if } Q_i \neq 0 \wedge M_{k,i} \neq 0
        \wedge Q_i = M_{k,i} \\
    0 & \text{otherwise}
  \end{cases}
\end{equation}

\subsection{The Role of the Null Trit}

This scoring function treats the null trit ($0$) differently
from $+1$ and $-1$:

\begin{itemize}[nosep]
  \item A null trit in the query means ``I don't have information
        about this dimension.'' It neither contributes to nor
        penalizes the score.
  \item A null trit in the motif means ``the perception system
        was uncertain about this dimension when the motif was
        committed.'' It neither contributes to nor penalizes the
        score.
  \item Only trits where \textit{both} query and motif are
        committed ($\pm 1$) and \textit{agree} contribute to
        the score.
\end{itemize}

This is fundamentally different from binary CAM matching, where
every bit contributes. In Krimelack, the null trit is \textit{not}
a value --- it is the absence of commitment. Two patterns can
match perfectly on all committed dimensions while having different
numbers of null trits. This is correct behavior: if the perception
system was uncertain about a dimension, that dimension should not
affect recognition.

\subsection{Best-Match Selection}

The motif with the highest score is returned as the best match.
Ties are broken by index (earliest stored motif wins --- temporal
priority).

\begin{equation}
  k^* = \arg\max_{k:\, k < D,\, \text{count}[k] > 0}
  \text{score}(Q, M_k)
\end{equation}

The recall outputs are:
\begin{itemize}[nosep]
  \item \texttt{best\_match}: the full motif $M_{k^*}$
  \item \texttt{match\_score}: $\text{score}(Q, M_{k^*})$
        as an 8-bit integer
  \item \texttt{recall\_valid}: $1$ if any match was found,
        $0$ if memory is empty
\end{itemize}

\subsection{Combinational Implementation}

The entire recall engine is implemented as combinational logic.
No clock is required. The query propagates through the comparator
array and the best-match result appears at the output wires
within one propagation delay.

For a 32-entry, 48-bit implementation:
\begin{itemize}[nosep]
  \item 32 parallel comparators ($\times$ 24 trits each)
  \item 32 score accumulators (8-bit each)
  \item 5-stage binary max-tree to find the best score
  \item Total: $\sim$2,000 equivalent gates
\end{itemize}

Recall latency is determined by propagation delay through the
comparator tree, approximately 3--5~ns in 7nm FinFET. This
matches the SPPU settling time, allowing the recall result to
feed back into the same perception cycle.

% ============================================================
\section{Clockless Habituation}

\subsection{The Corner-Trapping Problem}

A combinational perception system with a fixed dead zone
produces identical output for identical input. If the system
encounters a corner (a sensor state that produces a strong
coupling response), it will commit to the same decision
indefinitely. There is no escape. No clock-based timeout.
No iteration counter. The system is trapped by its own
correct behavior.

Traditional solutions require clocks or timers. Krimelack
solves this without either.

\subsection{Dead-Zone Feedback Mechanism}

The match score from recall feeds directly into the SPPU's
dead-zone threshold:

\begin{equation}
  \boxed{
  \tau_{\text{eff}} = \tau_{\text{base}} + f(\text{match\_score})
  }
\end{equation}

where $f$ is a damped mapping (score changes must exceed a
barrier of $\sim$15 units to propagate, preventing oscillation
from small score fluctuations).

\textbf{Familiar states} (high match score) raise the
commitment barrier. Trits that barely exceeded the dead zone
now fall below it and revert to null. The loom state changes.
A different decision emerges.

\textbf{Novel states} (low match score) leave the dead zone
at its baseline. Full sensitivity. Maximum coupling.

\textbf{No clock. No timer. No iteration.} The feedback is a
combinational loop: SPPU output $\to$ Krimelack recall $\to$
match score $\to$ SPPU dead zone $\to$ different SPPU output.
The loop settles to a fixed point where the dead-zone
adjustment is self-consistent.

\subsection{Cognitive Consequences}

The dead-zone feedback produces emergent behaviors that mirror
biological cognition:

\begin{center}
\begin{tabular}{@{}lp{9cm}@{}}
\toprule
\textbf{Behavior} & \textbf{Mechanism} \\
\midrule
\textbf{Novelty seeking} & Familiar states are harder to commit
  to; system gravitates toward unfamiliar states \\
\textbf{Attention} & Novel stimuli get full sensitivity; repeated
  stimuli get diminished sensitivity \\
\textbf{Boredom} & Sustained identical input raises dead zone
  until no trits commit; system goes silent \\
\textbf{Creativity} & Familiar responses suppressed; only novel
  decisions remain available \\
\textbf{Recognition} & Partial cue retrieves full motif; system
  ``remembers'' the structural state \\
\bottomrule
\end{tabular}
\end{center}

These behaviors are not designed or programmed. They are
mathematical consequences of dead-zone feedback from a
content-addressable memory. They emerge from the architecture.

\subsection{Stability Analysis}

The feedback loop is stable because:
\begin{enumerate}[nosep]
  \item The damping barrier prevents small fluctuations from
        propagating (hysteresis)
  \item The dead-zone increase is monotonic with match score
        (no oscillation)
  \item The maximum dead-zone increase is bounded (cannot
        exceed the maximum coupling pressure, which would
        prevent all commitment)
  \item The loop settles in one propagation delay (no
        multi-cycle ringing)
\end{enumerate}

% ============================================================
\section{TSAC Native Habituation (Future)}

On ferroelectric substrates (e.g., HfO$_2$-based Ternary
State Analog Cells), habituation becomes a free physical
property of the material.

\subsection{Mechanism}

Sustained polarization in one direction causes dielectric
fatigue --- the potential barrier rises with repeated identical
state:
\begin{itemize}[nosep]
  \item Sustained $+1$ $\to$ barrier: 20 $\to$ 22 $\to$ 25
        $\to$ 28
  \item Barrier exceeds coupling pressure $\to$ trit falls to null
  \item Null (rest) $\to$ barrier recovers $\to$ fresh sensitivity
\end{itemize}

The RC time constant of the ferroelectric material IS the
habituation rate. No circuit. No logic. Physics.

\subsection{Two-Layer Architecture}

On ASIC, both layers operate simultaneously:
\begin{itemize}[nosep]
  \item \textbf{Substrate layer} (TSAC fatigue): fast, local,
        per-trit. Physical depletion of the commitment mechanism.
  \item \textbf{Cognitive layer} (Krimelack feedback): slower,
        global, pattern-level. Experiential familiarity across
        the full loom state.
\end{itemize}

This mirrors biology: receptor saturation (fast, local) and
cortical habituation (slow, global) operating at different
scales.

% ============================================================
\section{Gate Count and Physical Implementation}

\subsection{CMOS Gate Count}

\begin{center}
\begin{tabular}{@{}lr@{}}
\toprule
\textbf{Component} & \textbf{Equivalent Gates} \\
\midrule
Motif register array (32 entries $\times$ 48 bits) & 1,536 \\
Hash table (32 entries $\times$ 8 bits) & 256 \\
Duplicate detection (32 parallel hash + exact compare) & 200 \\
Commit control (eligibility logic + write pointer) & 100 \\
Recall comparators (32 $\times$ 24-trit parallel match) & 1,600 \\
Score accumulators (32 $\times$ 8-bit) & 400 \\
Best-match tree (5-stage binary max) & 200 \\
Habituation feedback (match score $\to$ dead-zone adder) & 144 \\
\midrule
\textbf{Total per Krimelack block} & \textbf{$\sim$4,600} \\
\bottomrule
\end{tabular}
\end{center}

\subsection{Scaling}

\begin{center}
\begin{tabular}{@{}lrrr@{}}
\toprule
\textbf{Configuration} & \textbf{Entries} & \textbf{Gates} &
\textbf{Recall (7nm)} \\
\midrule
Minimal (FPGA demo) & 32 & 4,600 & $\sim$3~ns \\
Standard (ArcLoom-256) & 64 & 8,400 & $\sim$4~ns \\
Extended (vision) & 256 & 30,000 & $\sim$6~ns \\
\bottomrule
\end{tabular}
\end{center}

Recall latency scales logarithmically with depth (max-tree
depth = $\log_2 D$). Even at 256 entries, recall completes
within the SPPU settling window.

\subsection{Power}

Krimelack's recall path is combinational: zero dynamic power
between queries. The commit path requires a clock edge for
register writes, but commits are rare (only when L6 fires
and all eligibility criteria are met). Typical commit rate:
$<$1~Hz. Power contribution: negligible.

% ============================================================
\section{Comparison with Prior Art}

\subsection{Ternary Content-Addressable Memory (TCAM)}

TCAMs are deployed in every enterprise network router. They
match against three states (0, 1, X = don't care) in parallel.

\begin{center}
\begin{tabular}{@{}lll@{}}
\toprule
\textbf{Property} & \textbf{TCAM} & \textbf{Krimelack} \\
\midrule
Third state meaning & Wildcard (don't care) & Structural uncertainty \\
Match scoring & Exact match or fail & Graded (count of agreements) \\
Null contribution & Matches anything & Contributes nothing \\
Habituation & None & Dead-zone feedback \\
Commit gating & None & Quality-filtered \\
Application & Address lookup & Structural perception \\
\bottomrule
\end{tabular}
\end{center}

\subsection{Hopfield Networks}

Hopfield networks store patterns as energy minima and recall
via iterative energy minimization.

\begin{center}
\begin{tabular}{@{}lll@{}}
\toprule
\textbf{Property} & \textbf{Hopfield} & \textbf{Krimelack} \\
\midrule
Storage method & Distributed in weights & Explicit motif registers \\
Recall method & Iterative settling & Single combinational pass \\
Weight symmetry & Required ($W_{ij} = W_{ji}$) & Not applicable \\
Training & Hebbian learning rule & None (commit gating) \\
Capacity & $\sim$0.14$N$ patterns & $D$ entries (configurable) \\
Spurious attractors & Yes (mixture states) & No (exact storage) \\
Recall time & Multiple iterations & $\sim$3--5~ns \\
\bottomrule
\end{tabular}
\end{center}

\subsection{Memristor Habituation (Nature Communications, 2025)}

Third-order HfO$_2$/TiO$_x$ memristors with physical habituation.
Robot arm ignores 71\% of familiar stimuli.

\begin{center}
\begin{tabular}{@{}lll@{}}
\toprule
\textbf{Property} & \textbf{Memristor} & \textbf{Krimelack} \\
\midrule
Habituation level & Per-synapse (local) & Pattern-level (global) \\
Habituation type & Analog (continuous) & Digital (threshold) \\
Context & Embedded in trained SNN & Standalone, no ML \\
Recall type & SNN inference & Combinational CAM \\
Training required & Yes (SNN weights) & No \\
\bottomrule
\end{tabular}
\end{center}

\subsection{Summary: What Krimelack Has That Nobody Else Has}

\begin{enumerate}[nosep]
  \item \textbf{Structural motif storage} --- coupled multi-field
        ternary patterns, not bit vectors or weight distributions
  \item \textbf{Null-trit-aware scoring} --- uncertainty dimensions
        excluded from matching, not wildcarded
  \item \textbf{Quality-gated commits} --- dimensional exhaustion
        + uncertainty + resonance filtering prevents noise storage
  \item \textbf{Clockless pattern-level habituation} --- dead-zone
        feedback from match score, no clock, no timer
  \item \textbf{Combined in one architecture} --- no prior system
        provides all four
\end{enumerate}

% ============================================================
\section{Validation}

\subsection{FPGA Implementation}

Krimelack is implemented in Verilog
(\texttt{arcloom\_krimelack.v}, 167 lines) and validated on
Xilinx Zynq-7020 (PYNQ-Z2 board).

\subsection{Demonstrated Capabilities}

\begin{center}
\begin{tabular}{@{}lll@{}}
\toprule
\textbf{Capability} & \textbf{Status} & \textbf{Evidence} \\
\midrule
Motif commit gating (L6) & \textbf{Proven} & Motifs stored only on
  structural lock \\
Forced commit (AXI) & \textbf{Proven} & Software-triggered capture
  via register write \\
Duplicate detection & \textbf{Proven} & Same stimulus never produces
  two motifs \\
Cue-based recall & \textbf{Proven} & Partial query retrieves best
  match \\
Dead-zone feedback & \textbf{Proven} & Familiarity damping in live
  sensor loop \\
Autonomous navigation & \textbf{Proven} & 351 samples, 176s,
  3 IR + camera \\
Vision discrimination & \textbf{Proven (sim)} & 5 objects,
  0.11--0.36 match gaps \\
One-shot learning & \textbf{Proven} & Single turntable capture
  $\to$ recognition \\
\bottomrule
\end{tabular}
\end{center}

\subsection{Vision Results}

Five realistic objects discriminated using camera-derived motifs:

\begin{itemize}[nosep]
  \item Match gap between objects: 0.11--0.36 (robust discrimination)
  \item Same object with different noise: 0.895--0.954 (reliable recall)
  \item Memory per object: $\sim$4~KB structural fingerprint
  \item Training: zero. Single exposure via forced commit.
  \item Hardware: $\sim$4,600 gates of Krimelack + $\sim$2,250 gates
        of SPPU = $<$7,000 gates total
\end{itemize}

% ============================================================
\section{Applications}

\subsection{Immediate (FPGA)}

\begin{itemize}[nosep]
  \item Autonomous robot obstacle memory (``I've been here before'')
  \item Object recognition from turntable capture (one-shot learning)
  \item Sensor anomaly detection (novel states vs familiar baseline)
\end{itemize}

\subsection{Near-Term (ASIC)}

\begin{itemize}[nosep]
  \item Implantable medical sensors: cardiac rhythm motifs stored
        on-chip, anomaly = novel pattern not matching any stored motif
  \item Environmental monitoring: decade-life sensors that learn
        their ``normal'' environment and alert on structural novelty
  \item Safety-critical systems: deterministic pattern recognition
        with provable behavior (combinational, no software)
\end{itemize}

\subsection{Long-Term (TSAC)}

\begin{itemize}[nosep]
  \item Ferroelectric substrate: per-trit habituation as a free
        physical property
  \item Two-layer habituation: substrate (fast, local) + Krimelack
        (slow, global)
  \item Brain-scale structural memory: millions of motifs with
        physics-based adaptation
\end{itemize}

% ============================================================
\section{Conclusion}

Krimelack is a content-addressable structural memory architecture
that stores balanced ternary motifs with quality-gated commits,
parallel cue-based recall, and clockless pattern-level habituation.
It is the first memory system to combine all four capabilities in
a single architecture.

The null trit is not a wildcard or a don't-care bit. It is
structural uncertainty --- a first-class computational state that
influences matching, habituation, and dimensional exhaustion
detection. This distinction is fundamental: Krimelack does not
\textit{match patterns}. It recognizes structural states by
measuring how many committed dimensions agree. Dimensions where
the system was uncertain are honestly excluded.

The habituation mechanism --- dead-zone feedback from match
score --- produces novelty-seeking, attention, boredom, and
creativity as emergent consequences of the architecture, not as
designed behaviors. These cognitive properties arise from the
mathematics of a combinational feedback loop between a
content-addressable memory and a coupling fabric with dead zones.

Krimelack has been validated on FPGA silicon in autonomous robot
navigation with camera and IR sensor fusion. It operates at
$\sim$4,600 gates with combinational recall latency of
$\sim$3--5~ns. No prior art exists for the combined architecture.

\vspace{2em}
\hrule
\vspace{0.5em}
{\color{nm-gray}\small This document is confidential and for
internal use only. Krimelack is a component of the ArcLoom SPPU
architecture, the intellectual property of Joseph Forrester /
DSF-AI. The coupling weight derivation methodology and UF pipeline
internals are trade secrets and are not disclosed in this document.}

\end{document}
